\section{Macchine di Turing}
\subsection{Definizione}
\begin{itemize}
    \item $M = (Q, \Sigma, \delta, s)$: Macchina di Turing deterministica.
    \item $N = (Q, \Sigma, \Delta, s)$: Macchina di Turing non deterministica
    \item $k \in \mathbb{N}$: Numero di nastri.
\end{itemize}

\hrulefill

\begin{itemize}
    \item $Q$: Insieme finito di stati.
    \item $s \in Q$: Stato iniziale.
    \item $H = \{h, "yes", "no"\} \subset Q$: Insieme di stati terminali.
\end{itemize}

\hrulefill

\begin{itemize}
    \item $\Sigma$: Alfabeto (insieme finito di simboli) con cui viene codificato l'input.
    \item $\triangleright$: Start.
    \item $\sqcup$: Blank.
    \item $\Sigma_\sqcup = \Sigma \cup \{\sqcup\}$
    \item $\Gamma = \Sigma_\sqcup \cup \{\triangleright\}$ Alfabeto di lavoro che la macchina utilizza per letture e scritture.
\end{itemize}

\hrulefill

\begin{itemize}
    \item $\{\leftarrow, \rightarrow, -\}$ spostamenti della testina sul nastro: puó muoversi di una cella a sinistra o a destra, oppure restare ferma.
    \item $\delta: Q \times \Gamma^k \rightarrow Q \times \Gamma^k \times \{\leftarrow, \rightarrow, -\}^k$ \underline{funzione} di transizione: prende in input i simboli sotto le testine e lo stato corrente, restituisce simboli che sovrascrivono quelli letti, uno stato in cui transitare al passo successivo e gli spostamenti che ogni testina deve effettuare effettuare alla fine.
    \begin{itemize}
        \item $\delta(q, [\dots, \triangleright_i, \dots]) = (q', [\dots, \triangleright_i, \dots], [\dots, \rightarrow_i, \dots])$ (con $q, q' \in Q$): se una qualsiasi testina $i$ si trova su $\triangleright$, deve riscriverlo e spostarsi obbligatoriamente a destra.
    \end{itemize}
    \item $\Delta \subseteq (Q \times \Gamma^k) \times (Q \times \Gamma^k \times \{\leftarrow, \rightarrow, -\}^k)$ \underline{relazione} di transizione: associa a ogni coppia di stato corrente e simboli sotto la testina, una o piú terne di nuovi stati, simboli da scrivere e direzioni verso cui spostare le testine.
    \begin{itemize} 
        \item $C_{(q,\sigma)} = \{(q', \sigma', \textit{dir}) \mid (q, \sigma, \quad q', \sigma', \textit{dir}) \in \Delta\}$ insieme delle \textit{scelte} che $N$ puó fare a partire dal \textit{contesto} $(q, \sigma)$.
        \item $d = \max_{(q, \sigma) \in Q \times \Gamma^k} |C_{(q,\sigma)}|$ grado di non determinismo di $N$.
        \subitem Ovviamente $d$ é sempre finito perché $\Delta$ é un prodotto di insiemi finiti.
    \end{itemize}
\end{itemize}

\hrulefill

\begin{itemize}
    \item $x \in \Sigma^*$: Input scritto sul primo nastro all'inizio della computazione.
    \item $n = |x| \in \mathbb{N}$
    \item $f(n)$: Tempo utilizzato dalla computazione.
    \item $g(n)$: Spazio di memoria utilizzato dalla computazione.
\end{itemize}


\subsubsection{Non determinismo}
La differenza tra una macchina di Turing deterministica e una non deterministica 
é che la prima, a ogni passo, ha una sola scelta possibile, mentre la seconda puó 
avere piú scelte possibili e seguirle in maniera simultanea. \\
Vedremo che la computazione di una MdT non deterministica puó essere rappresentata 
come un albero. \\
Il non determinismo non va confuso con il parallelismo, che é rappresentato dai $k$ nastri. 
Il non determinismo non consiste in diverse unitá di computazione che eseguono 
contemporaneamente istruzioni diverse, ma in diverse scelte che vengono fatte in maniera 
simultanea e senza problemi di sincronizzazione. \\
Un altro modo di vederlo é con l'idea di un oracolo, che "suggerisce" alla macchina la strada
giusta da seguire per trovare la soluzione al problema. \\
Ma forse il modo piú concreto é quello di vedere gli algoritmi non deterministici come 
algoritmi che si limitano a controllare che una soluzione risolva il problema, senza 
preoccuparsi di cercarla. \\
Molti problemi infatti richiedono, per trovare la soluzione ottima, di esplorare uno spazio 
di soluzioni molto grande, spesso esponenziale rispetto alla dimensione dell'input. \\
Le macchine non deterministiche scorporano il costo richiesto dall'esplorazione delle soluzioni.


\subsection{Configurazione}
Siano, $\forall i \in \{1, \dots, k\}$:
\begin{itemize}
    \item $q \in Q$: Stato corrente.
    \item $w_i = "\triangleright w_{i,2} w_{i,3} \dots {w_{i,{m_i}}}" \in \{\triangleright\} \times {\Sigma_\sqcup}^*$: 
    Stringa sul nastro $i$ prima della testina, incluso il simbolo sotto la testina, $w_{i,m_i}$. La sua lunghezza, $m_i$, puó essere diversa da nastro a nastro.
    \item $u_i = "u_{i,1} u_{i,2} \dots u_{i,{p_i}}" \in {\Sigma_\sqcup}^*$: Stringa sul nastro $i$ dopo la testina. Come prima, la sua lunghezza, $p_i$, puó essere diversa da nastro a nastro.
\end{itemize}

\[
\configurazione{q}{w}{u}
\]

\subsubsection{Massimo numero di configurazioni possibili}
\begin{itemize}
    \item Numero di stati: $|Q|$
    \item Massimo numero di stringhe possibili: $|\Sigma_\sqcup|^{g(n) \cdot k}$ \\
    ($|\Sigma_\sqcup|^{g(n)}$ stringhe per ogni nastro)
    \item Massimo numero di posizioni delle testine possibili: $g(n)^k$ \\
    ($g(n)$ posizioni per ogni testina)
\end{itemize}

{\[|Q| \cdot |{\Sigma_\sqcup|}^{g(n) \cdot k} \cdot g(n)^k\]}


\subsection{Computazione}
\subsubsection{Computazione di un singolo passo}
\[
\configurazione{q}{w}{u}
\xrightarrow{M}
\configurazione{q'}{w'}{u'}
\]
Se e solo se, dato:
\begin{itemize}
\item MdT deterministiche: $\qquad \delta(q, [{w_m}]) = (q', [\sigma], [\textit{dir}])$.
\item MdT non deterministiche: $(q, [w_m],\quad q', [\sigma], [\textit{dir}]) \in \Delta$.
\item Dove:
\begin{itemize}
	\item $[{w_m}] = ({w_{1,m_1}}, \dots, {w_{k,m_k}})$: Simboli sotto le testine dei $k$ nastri.
    \item $q' \in Q$: Stato successivo della macchina.
    \item $[\sigma] = [\sigma_1, \dots, \sigma_k] \in {\Sigma_\sqcup}^k$: Vettore di $k$ caratteri, ognuno da scrivere su una delle $k$ celle sotto una testina.
    \item $[\textit{dir}] = [\textit{dir}_1, \dots, \textit{dir}_k] \in \{\leftarrow, \rightarrow, -\}^k$: Vettore di $k$ direzioni che ogni testina deve seguire.
\end{itemize}
\end{itemize}
Vale che, $\forall i \in \{1, \dots, k\} :$
{\small
\begin{itemize}
    \item $(\textit{dir}_i = -) \implies ((w'_i = "\triangleright w_{i,2} w_{i,3} \dots {\sigma_i}") \land (u'_i = u_i))$
    \item $(\textit{dir}_i = \leftarrow) \implies ((w'_i = "\triangleright w_{i,2} w_{i,3} \dots {w_{i,m_{i-1}}}") \land (u'_i = "\sigma_i u_{i,1} u_{i,2} \dots u_{i,p_i}"))$
    \item $((\textit{dir}_i = \rightarrow) \land (u_i \neq \epsilon)) \implies ((w'_i = "\triangleright w_{i,2} w_{i,3} \dots \sigma_i {u_{i,1}}") \land (u'_i = "u_{i,2} \dots u_{i,p_i}")) $
    \item $((\textit{dir}_i = \rightarrow) \land (u_i = \epsilon)) \implies ((w'_i = "\triangleright w_{i,2} w_{i,3} \dots \sigma_i {\sqcup}") \land (u'_i = \epsilon)) $
\end{itemize}
}

\subsubsection{Computazione di $t$ passi}
\[
\configurazione{q_0}{w_0}{u_0}
\xrightarrow{M^t}
\configurazione{q_t}{w_t}{u_t}
\]

Se $\forall j \in \{0, \dots, t-1\} :$
\[
\configurazione{q_j}{w_j}{u_j}
\xrightarrow{M}
\configurazione{q_{j+1}}{w_{j+1}}{u_{j+1}}
\]

\subsubsection{Computazione di un numero di passi indefinito}
\[
\configurazione{q}{w}{u}
\xrightarrow{M^*}
\configurazione{q'}{w'}{u'}
\]

Se $\exists t \in \mathbb{N} :$
\[
\configurazione{q}{w}{u}
\xrightarrow{M^t}
\configurazione{q'}{w'}{u'}
\]

\subsection{Terminazione}
    $M(x) \in \Sigma^* \cup \{"yes", "no",\nearrow$\}
    \begin{itemize}
        \item $M(x) = "yes"$ sse $\exists w, u :$
        \[
        \configurazionestart{x} \xrightarrow{M^*} \configurazione{"yes"}{w}{u}
        \]
        \item $M(x) = "no"$ sse $\exists w, u :$
        \[
        \configurazionestart{x} \xrightarrow{M^*} \configurazione{"no"}{w}{u}
        \]
        \item $M(x) = y \in \Sigma^*$ sse $\exists w, u :$
        \[
        \configurazionestart{x} \xrightarrow{M^*} \configurazione{h}{w}{u} \qquad \land \qquad 
        y = "w_k u_k"
        \]
        \item $M(x) = \nearrow$ se la macchina diverge (la computazione non termina). Cioé se non vale nessuno dei casi precedenti.
	\end{itemize}

\subsection{Macchine di Turing con input e output}
\begin{itemize}
	\item $k \geq 2$
	\item Il primo nastro non viene mai modificato (read only).
	\item L'ultimo nastro viene scandito solo verso destra (write only).
	\item Il cursore del primo nastro non supera mai il primo $\sqcup$ successivo all'input $x$.
\end{itemize}

\subsection{Composizione di Macchine di Turing}
Date due MdT $M_1 = (Q_1, \Sigma, \delta_1, s_1)$ e $M_2 = (Q_2, \Sigma, \delta_2, s_2)$ tali che:
\begin{itemize}
	\item Hanno lo stesso alfabeto $\Sigma$ 
	\subitem (Se gli alfabeti sono diversi, basta considerare la loro unione).
	\item $Q_1 \cap Q_2 = \emptyset$ 
	\subitem (Se gli insiemi di stati non sono disgiunti, basta rinominare gli stati di uno dei due).
	\item $M_1$ termina con tutte le testine su $\triangleright$ 
	\subitem (Se non accade, basta cambiare la funzione di transizione $\delta_1$, redirezionando le coppie che portano verso gli stati $"yes"$ e $h$ in uno stato che riporta le testine su $\triangleright$ e poi termina).
\end{itemize}

La loro \textbf{composizione} é: $(M_1, M_2) =  (Q, \Sigma, \delta, s)$, dove: 
\begin{itemize}
	\item $Q = Q_1 \cup Q_2$
	\item $s = s_1$
	\item $\forall \sigma \in \Sigma^k$:
	\[
	\delta(q, \sigma) =
	\begin{cases} 
		\delta_2(q, \sigma)  \quad se \quad q \in Q_2 \\
		\delta_1(q, \sigma) \quad se \quad q \in Q_1 \land \delta_1(q, \sigma)  = (p, \tau, \textit{dir}) \land p \notin \{h, "yes"\} \\
		(s_2, \tau, \textit{dir}) \quad se \quad q \in Q_1 \land  \delta_1(q, \sigma)  = (p, \tau, \textit{dir}) \land p \in \{h, "yes"\} \\
	\end{cases}
	\]
	Dove $\tau \in \Sigma^k$ \quad $\textit{dir} \in \{\leftarrow, \rightarrow, -\}^k$
\end{itemize}

\begin{itemize}
	\item La funzione di transizione modifica solo $\delta_1$, reindirizzando le coppie $(q, \sigma)$ che portano a uno stato terminale $h$ o $"yes"$ verso lo stato iniziale di $M_2$. 
	\item Tutti gli altri casi restano uguali alle funzioni di transizione delle due macchine.
	\item Le coppie $(q, \sigma)$ che per $\delta_1$ vanno nello stato $"no"$ non devono essere reindirizzate: \textbf{se $\mathbf{M_1}$ rigetta l'input, la computazione non deve continuare}.
\end{itemize}

% \newpage
\subsection{Esempi}
Gli esempi possono essere testati su \url{https://morphett.info/turing/turing.html}

\subsubsection{Traslazione di una stringa binaria}
\label{example:shift}
\[
\begin{array}{c|c|c|c|c}
\text{Stato} & \text{Letto} & \text{Scritto} & \text{Movimento} & \text{Nuovo stato} \\ \hline
s   & \triangleright & \triangleright & \rightarrow & s \\
s   & \sqcup         & \sqcup         & \leftarrow  & q \\
s   & *              & *              & \rightarrow & s \\
q   & \triangleright & \triangleright & \rightarrow & h \\
q   & 1              & \sqcup         & \rightarrow & q_1 \\
q   & 0              & \sqcup         & \rightarrow & q_0 \\
q_1 & \sqcup         & 1              & \leftarrow  & s \\
q_0 & \sqcup         & 0              & \leftarrow  & s
\end{array}
\]
